% Source for KEYS.pdf, the one-page printable key reference.
% Build:  pdflatex -interaction=nonstopmode KEYS.tex
\documentclass[10pt]{article}
\usepackage[letterpaper,margin=0.5in]{geometry}
\usepackage[T1]{fontenc}
\usepackage[utf8]{inputenc}
\usepackage{array}
\usepackage{textcomp}
\pagestyle{empty}
\setlength{\parindent}{0pt}
\setlength{\tabcolsep}{5pt}
\renewcommand{\arraystretch}{1.22}

% Key names are set in monospace; arrows come from math mode.
\newcommand{\key}[1]{\texttt{#1}}
\newcommand{\up}{$\uparrow$}
\newcommand{\dn}{$\downarrow$}
\newcommand{\lft}{$\leftarrow$}
\newcommand{\rgt}{$\rightarrow$}
\newcommand{\sep}{\hspace{0.7em}}
\newcommand{\heading}[1]{\par\vspace{5pt}\textbf{#1}\par\vspace{3pt}}

% A two-column key table with a rule under the header and between every row.
\newenvironment{keytable}[1]
  {\begin{tabular}{|>{\raggedright\arraybackslash}p{5.4cm}|>{\raggedright\arraybackslash}p{#1}|}\hline
   \textbf{Key} & \textbf{Action}\\\hline\hline}
  {\end{tabular}}

\begin{document}
\footnotesize

{\large\textbf{beeswax --- key reference}}\\[3pt]
Navigation is vi-style; arrow keys, Home/End/PgUp/PgDn and the F-keys work
everywhere too, so nothing requires the vi bindings.\\[3pt]
\textbf{Unprefixed} --- vi verbs and motions, acting on the item under the cursor.
\textbf{Shift} --- the reverse, or the larger version, of the unshifted key.
\textbf{Ctrl} --- whole-file and whole-view commands, and moving structure.
\textbf{Alt} --- adding structure around the items: sections and columns.\\[3pt]
In the Key column, \key{/} separates keys matching the parts of the Action in
order; a wider gap separates keys that do the same thing.\\[3pt]
Structural editing follows VimOutliner, which leaves it to plain vim:
\key{o}/\key{O} makes a sibling, \key{>{}>}/\key{<{}<} changes its level, \key{u}
and \key{Ctrl+R} undo and redo.

\heading{Moving around}
\begin{keytable}{12.6cm}
\key{k} / \key{j} \sep \up\ / \dn & Up / down\\\hline
\key{h} / \key{l} \sep \lft\ / \rgt \sep \key{Tab} / \key{Shift+Tab} & Left / right across columns\\\hline
\key{Enter} / \key{Shift+Enter} & Down / up\\\hline
\key{PgDn} / \key{PgUp} \sep \key{Ctrl+F} / \key{Ctrl+B} & One screenful down / up\\\hline
\key{Home} / \key{End} \sep \key{\{} / \key{\}} & Section head / last item; again for the previous or next section\\\hline
\key{Shift+Home} / \key{Shift+End} & Head of the previous / next section\\\hline
\key{Ctrl+Home} / \key{Ctrl+End} \sep \key{gg} / \key{G} & Top / bottom of the view\\\hline
\key{H} / \key{M} / \key{L} & First / middle / last visible row\\\hline
\key{zz} / \key{z Enter} / \key{z-} & Put the cursor at center / top / bottom\\\hline
\key{Ctrl+PgUp} / \key{Ctrl+PgDn} & Previous / next view\\\hline
\key{/} & Open the search bar, forwards\\\hline
\key{?} & Open the search bar, backwards\\\hline
\key{n} / \key{N} & Repeat the search, same / opposite direction (wraps)\\\hline
\end{keytable}

\heading{Items}
\begin{keytable}{12.6cm}
\key{i} \sep \key{F2} & Edit item text or cell value\\\hline
\key{o} / \key{O} \sep \key{Insert} & New item below / above\\\hline
\key{x} \sep \key{Delete} & Remove from this section's category; on a head remove the section or column; in a data column clear the value\\\hline
\key{Shift+Delete} & Discard from all categories, permanently\\\hline
\key{u} / \key{Ctrl+R} \sep \key{Ctrl+Z} / \key{Ctrl+Y} & Undo / redo\\\hline
\key{F3} / \key{F4} / \key{F5} & Assign categories or pick a value / mark done / note in \texttt{\$EDITOR}\\\hline
\key{F6} & Item, section or column properties\\\hline
\key{F8} / \key{F9} / \key{F10} & View Manager / Category Manager / Menu\\\hline
\key{Ctrl+S} / \key{Alt+Q} / \key{F1} & Save / quit / help\\\hline
\end{keytable}

\heading{Structure --- View screen}
\begin{keytable}{12.6cm}
\key{Alt+j} / \key{Alt+k} \sep \key{Alt+}\dn\ / \key{Alt+}\up & New section below / above\\\hline
\key{Alt+h} / \key{Alt+l} \sep \key{Alt+}\lft\ / \key{Alt+}\rgt & New column left / right\\\hline
\key{Ctrl+h} / \key{Ctrl+l} \sep \key{Ctrl+}\lft\ / \key{Ctrl+}\rgt & Move this column left / right\\\hline
\key{Alt+s} & Sort this section now\\\hline
\end{keytable}

Alt plus a direction adds; Ctrl plus a direction moves. The adds are quick: a
category picker opens and the section or column appears as soon as a category
is chosen. The full Section Add and Column Add forms are on the F10 menu.

\heading{Structure --- Category Manager (F9)}
\begin{keytable}{12.6cm}
\key{o} / \key{O} \sep \key{Insert} & New sibling below / above\\\hline
\key{>{}>} \sep \key{Ctrl+l} \sep \key{Ctrl+}\rgt & Demote one level --- \key{o} then \key{>{}>} makes a child\\\hline
\key{<{}<} \sep \key{Ctrl+h} \sep \key{Ctrl+}\lft & Promote one level\\\hline
\key{Ctrl+j} / \key{Ctrl+k} \sep \key{Ctrl+}\dn\ / \key{Ctrl+}\up & Move down / up among siblings\\\hline
\key{Enter} \sep \key{F2} \sep \key{i} & Rename\\\hline
\key{Alt+F10} & Move mode: plain arrows, Enter when done\\\hline
\key{Delete} \sep \key{x} & Delete\\\hline
\key{F6} & Category properties\\\hline
\end{keytable}

On a root category \key{o} and \key{O} both create a first child, a root having
no siblings. \key{Alt+F10} is the only Alt key on this screen.

\heading{Lists and pickers}
The Category Manager, View Manager, F3 pickers, Section Select and the sort and
filter dialogs share one set: \key{k}/\key{j}, \key{PgUp}/\key{PgDn},
\key{gg}/\key{G}, \key{H}/\key{M}/\key{L}, \key{/} to search,
\key{i}/\key{o}/\key{x}/\key{u} to rename, create, delete and undo where
supported, \key{Enter} to confirm, \key{Esc} to cancel. In the View Manager and
in View Properties' Sections field, \key{Ctrl+j} / \key{Ctrl+k} reorder the
list, as they do in the Category Manager.\\[3pt]
The category pickers --- the quick adds, and \key{F3} on the Section Add form
--- manage categories in place: \key{/} searches (\key{F7} / \key{F8} step through
matches), \key{Insert} adds, \key{F2} renames, \key{Delete} discards, \key{F6}
opens properties. \key{Esc} there closes the picker and returns to the dialog
rather than cancelling it. \key{Delete} is refused for a protected category and
for the current section's head category.

\end{document}
